\documentclass[11pt]{article}

\usepackage[utf8]{inputenc}
\usepackage[T1]{fontenc}
\usepackage[a4paper,margin=2.5cm]{geometry}
\usepackage{amsmath}
\usepackage{amssymb}
\usepackage{booktabs}
\usepackage{microtype}
\usepackage[colorlinks=true,linkcolor=blue,citecolor=blue,urlcolor=blue]{hyperref}

\title{Evaluation and Comparison Extensions to the Tile-Level Pruning Strategy}
\author{Sebastian Fiebiger}
\date{July 16, 2026}

\begin{document}

% ---------------- Title page ----------------
\begin{titlepage}
\centering
\vspace*{4cm}
{\Huge\bfseries Evaluation and Comparison Extensions\\[0.3em]
to the Tile-Level Pruning Strategy\par}
\vspace{1.1cm}
{\large Controls, Sparsity Sweeps, and Additional Evaluation Metrics for\\
Tile-Level Redundancy Detection in Large Language Models\par}
\vspace{2.2cm}
{\large\bfseries Sebastian Fiebiger\par}
\vspace{0.3cm}
{\normalsize Technical University of Darmstadt\par}
\vfill
{\normalsize July 16, 2026\par}
\end{titlepage}

% ---------------- Abstract + Contents ----------------
\begin{abstract}
This note records the extensions introduced after a review of the tile-level pruning strategy for
Wanda and SparseGPT. The staged pruning strategy is adopted without modification; the extensions
described here strengthen the experimental design that surrounds it rather than changing the pruning
procedures themselves. Four additions are proposed: (i)~random and magnitude tile pruning as explicit
controls; (ii)~repetition of random pruning across multiple seeds to quantify variance; (iii)~treatment
of tile sparsity as a primary experimental axis through a perplexity-versus-sparsity sweep; and
(iv)~evaluation metrics beyond perplexity, namely the output-distribution divergence between the dense
and pruned models and downstream task accuracy. For each addition the motivation, the method, and its
relationship to the adopted plan are stated, followed by a consolidated execution schedule with runtime
estimates. All experiments target Qwen3-4B with $32\times32$ tiles on the representative layers
$0, 9, 18, 27$ and $35$, consistent with the adopted strategy.
\end{abstract}

\tableofcontents
\newpage

% ---------------- Sections ----------------
\section{Introduction and Scope}

The tile-level redundancy project studies which tiles of a transformer's weight matrices can be removed
while preserving model behaviour. A staged pruning strategy based on Wanda and SparseGPT defines how
tiles are selected and at what scope: a single matrix, a complete layer, a cluster of layers, and finally
the complete model.

This document does not alter that strategy. Its purpose is to record the design extensions added after
review so that the resulting measurements are interpretable, statistically meaningful, and not dependent
on a single evaluation metric. The extensions fall into two groups: additional pruning methods that serve
as controls (Sections~\ref{sec:controls} and~\ref{sec:seeds}), and additional ways of measuring the
quality of a pruned model (Sections~\ref{sec:sparsity} and~\ref{sec:eval}). Section~\ref{sec:distinction}
first clarifies a distinction that motivates the rest.

\section{Pruning Metrics versus Evaluation Metrics}
\label{sec:distinction}

Two different quantities are both loosely called ``metrics'', and separating them is essential to the
design that follows.

\begin{itemize}
  \item A \textbf{pruning metric} decides \emph{which} tiles to remove. Random selection, weight
  magnitude, Wanda, and SparseGPT reconstruction error are pruning metrics.
  \item An \textbf{evaluation metric} decides \emph{whether} the resulting pruned model is still good.
  Perplexity, output-distribution divergence, and downstream task accuracy are evaluation metrics.
\end{itemize}

In particular, SparseGPT is a \emph{pruning} metric, not an evaluation metric: its reconstruction error
is a local, per-layer signal used while pruning, and it does not by itself indicate whether the complete
model still performs its tasks. The extensions in Sections~\ref{sec:sparsity} and~\ref{sec:eval}
therefore concern evaluation metrics, which were previously limited to perplexity.

\section{Extension 1: Baseline Pruning Methods as Controls}
\label{sec:controls}

Two reference methods are added alongside the data-aware methods: \textbf{random} tile pruning and
\textbf{magnitude} tile pruning.

\paragraph{Rationale.} The result of a data-aware method is only interpretable relative to a control. A
statement such as ``Wanda removes 20\% of tiles with a small perplexity increase'' has no meaning without
a floor (random selection) and a naive reference (weight magnitude). These two controls reinstate the
four-method comparison already recommended in the earlier metric report.

\paragraph{Requirement.} The controls must be evaluated at identical settings --- model, tile size,
sparsity, and evaluation data --- as the data-aware methods. They are therefore run within the same sweeps
rather than deferred to a separate later phase, because a control produced at mismatched settings is not
a valid comparison.

\section{Extension 2: Multiple Seeds and Variance}
\label{sec:seeds}

Random tile pruning is repeated across \textbf{five seeds}, and its results are reported as a mean together
with a range.

\paragraph{Rationale.} A conclusion such as ``layer 18 is robust'', drawn from a single run at a single
sparsity, may reflect noise rather than a real property of the model. Reporting the spread of the random
control makes the robust-versus-sensitive comparison honest and quantifies how large a perplexity
difference must be before it is meaningful.

The remaining methods --- magnitude, Wanda, and both SparseGPT variants --- are deterministic: repeated
runs reproduce their results exactly, which was confirmed empirically (bit-identical perplexity across
independent runs). A single run therefore suffices for each of them, and only random pruning requires
multiple seeds.

\section{Extension 3: Sparsity as a Primary Axis}
\label{sec:sparsity}

Rather than fixing a single sparsity, model quality is measured as a function of sparsity, at two
resolutions chosen according to what each stage requires. The screening and clustering stages, which locate
redundancy across many layer--matrix combinations, use a coarse grid
\[
  \{\,10,\ 20,\ 40\,\}\,\%,
\]
sufficient to reveal whether each combination degrades gradually or collapses sharply, and whether the
robust-versus-sensitive ranking is stable as pruning becomes more aggressive. The final whole-model
comparison, which requires a smooth curve, uses a fine grid
\[
  \{\,5,\ 10,\ 20,\ 30,\ 40,\ 50,\ 60,\ 70\,\}\,\%,
\]
and the resulting perplexity-versus-sparsity curve is treated as a primary result of the study.

\paragraph{Rationale.} Redundancy is sparsity-dependent: a matrix that appears robust at 20\% may collapse
at 40\%. A redundancy map built at a single sparsity can therefore mislabel a region, which is why even the
screening stages span several sparsities rather than one. A coarse grid is adequate there because the
purpose is to rank many combinations, not to draw a precise curve for each; the fine grid is reserved for
the whole-model curve, where the location of the collapse and the separation between methods must be
resolved exactly. Pruning methods are nearly indistinguishable at low sparsity and separate only as sparsity
increases, so the region where a data-aware method continues to hold while the random control has already
collapsed is the most informative part of the curve, and capping the sweep prematurely would hide exactly
where a method proves its value.

The whole-model sweep is extended until all methods, including the strongest, have clearly degraded. Because
perplexity grows by orders of magnitude near collapse, the perplexity axis is plotted on a logarithmic scale.

\section{Extension 4: Additional Evaluation Metrics}
\label{sec:eval}

Perplexity is retained as the primary language-modelling metric but is complemented by two further
measurements, so that no conclusion rests on a single number.

\subsection{Output-distribution divergence}

For the same inputs, the divergence between the dense and pruned models' output distributions is recorded.
The Kullback--Leibler divergence of the output distributions is
\[
  D_{\mathrm{KL}}\!\left(P_{\mathrm{dense}} \,\big\|\, P_{\mathrm{pruned}}\right)
  = \sum_{i} P_{\mathrm{dense}}(i)\,\log\frac{P_{\mathrm{dense}}(i)}{P_{\mathrm{pruned}}(i)} .
\]
This is reported together with the top-1 agreement rate (the fraction of positions at which both models
predict the same next token) and the cosine similarity of the hidden states,
\[
  \cos\!\left(h_{\mathrm{dense}}, h_{\mathrm{pruned}}\right)
  = \frac{\langle h_{\mathrm{dense}},\, h_{\mathrm{pruned}}\rangle}
         {\lVert h_{\mathrm{dense}}\rVert \, \lVert h_{\mathrm{pruned}}\rVert} .
\]
These quantities are inexpensive and highly sensitive, and are therefore used at the screening stage, where
they measure behavioural drift token by token.

\subsection{Downstream task accuracy}

For the final whole-model variants, accuracy on HellaSwag, PIQA and ARC-Easy is measured with the
lm-evaluation-harness, which is already configured and baselined for the dense model. This tests whether
the pruned model retains actual task ability.

\paragraph{Rationale.} A model may preserve a reasonable perplexity while quietly degrading on downstream
tasks. Evaluating along three complementary axes --- language-modelling loss, behavioural fidelity, and
task accuracy --- prevents the study's conclusions from resting on a single, potentially misleading
measurement.

\section{Consolidated Execution Plan}

The extensions are folded into the adopted execution order without changing its stages. Each stage
additionally runs the control methods and, where relevant, the sparsity grid. The estimated schedule on a
single RTX~4080~Super is summarised below.

\begin{table}[h]
\centering
\small
\begin{tabular}{@{}llccc@{}}
\toprule
Stage & Scope & Methods & Sparsity & Est.\ time \\
\midrule
0\ \ Dense baseline    & full model                  & dense reference & ---          & done \\
1\ \ Per-matrix screen & 5 layers $\times$ 7 matrices & all + controls  & $10/20/40\%$ & $\sim$6\,h \\
2\ \ Whole-layer       & 5 layers                    & all + controls  & $10/20/40\%$ & $\sim$2\,h \\
3\ \ Cluster analysis  & robust / sensitive clusters & key methods     & $10/20/40\%$ & $\sim$4\,h \\
4\ \ Sparsity sweep    & whole model                 & all + controls  & 8-point grid & $\sim$4\,h \\
5\ \ Final + downstream& whole model                 & best methods    & matched      & $\sim$3\,h \\
\bottomrule
\end{tabular}
\end{table}

The total is approximately nineteen hours of GPU time, well within the available budget. A fixed perplexity
subset is used during the screening stages, where measurements serve only to rank layer--matrix
combinations; full evaluation is reserved for the final and headline results, following the cost-control
guidance of the adopted strategy.

\section{Relationship to the Adopted Strategy}

The backbone is unchanged: the model (Qwen3-4B), the tile size ($32\times32$), the representative layers
($0, 9, 18, 27, 35$), the dense-model discipline of restoring the original weights between independent
experiments, the Wanda and SparseGPT strategy ladders, and the combined variant that selects tiles with
Wanda and reconstructs the survivors with SparseGPT.

Of the extensions, the controls of Sections~\ref{sec:controls} and~\ref{sec:seeds} reinstate elements
already present in the earlier metric report. The additional evaluation metrics of Section~\ref{sec:eval}
and the expanded sparsity treatment of Section~\ref{sec:sparsity} extend the strategy document, which named
perplexity as the initial evaluation metric and scoped sparsity sweeps narrowly. None of the extensions
contradicts the adopted strategy; each strengthens the evidence it produces.

\section{Conclusion}

These extensions convert a single-metric, single-sparsity procedure into a controlled, multi-sparsity,
multi-metric comparison. The pruning strategy is preserved in full. The additions ensure that its results
can be interpreted against controls, carry a sense of variance, are characterised across the full sparsity
range, and are corroborated by more than one measure of model quality.

\begin{thebibliography}{9}
\bibitem{wanda} M.~Sun, Z.~Liu, A.~Bair, and J.~Z.~Kolter, ``A simple and effective pruning approach for
large language models,'' in \emph{The Twelfth International Conference on Learning Representations}, 2024.
\bibitem{sparsegpt} E.~Frantar and D.~Alistarh, ``SparseGPT: Massive language models can be accurately
pruned in one-shot,'' in \emph{Proceedings of the 40th International Conference on Machine Learning},
PMLR, vol.~202, 2023, pp.~10323--10337.
\end{thebibliography}

\end{document}
